% Generated draft: requires mathematical and source review.
\documentclass[11pt]{article}
\usepackage[margin=25mm]{geometry}
\usepackage{fontspec}
\setmainfont{Latin Modern Roman}
\usepackage{amsmath,amssymb,mathtools,unicode-math}
\setmathfont{Latin Modern Math}
\usepackage{enumitem,needspace,xurl,hyperref}
\setlength{\parindent}{0pt}
\setlength{\parskip}{5pt}
\setlength{\emergencystretch}{3em}
\newcommand{\N}{\mathbb{N}}
\newcommand{\Z}{\mathbb{Z}}
\newcommand{\Q}{\mathbb{Q}}
\newcommand{\R}{\mathbb{R}}
\newcommand{\C}{\mathbb{C}}
\newcommand{\F}{\mathbb{F}}
\newcommand{\A}{\mathbb{A}}
\newcommand{\PP}{\mathbb{P}}
\newcommand{\E}{\mathbb{E}}
\newcommand{\eps}{\varepsilon}
\newcommand{\dd}{\,\mathrm{d}}
\begin{document}
\section*{UnsolvedMath ID: 30004529}
\textbf{Source problem code:} catalogue entry
\subsection{Definitions and mathematical statement}
The source does not specify the ground field or the parameter space underlying the word ``family.'' The following definitions describe a conditional interpretation, whose choices require confirmation.

Fix a field \(k\). Under the convention that the algebras are commutative, finite-dimensional local \(k\)-algebras \((A,\mathfrak m)\) with residue field \(A/\mathfrak m\cong k\), define
\[
\operatorname{Soc}(A)=\{a\in A:\mathfrak m a=0\},\qquad
h_A(i)=\dim_k(\mathfrak m^i/\mathfrak m^{i+1})\quad(i\geq0).
\]
Such an algebra is Gorenstein precisely when \(\dim_k\operatorname{Soc}(A)=1\). In this interpretation, a Gorenstein sequence is a finite sequence \(T=(t_0,\ldots,t_s)\), with \(t_s\ne0\), realized by such an algebra with
\[
h_A(i)=t_i\quad(0\leq i\leq s),\qquad h_A(i)=0\quad(i>s).
\]
These are local algebras; no grading or symmetry of \(T\) is imposed here.

Write \(\operatorname{Gor}(T)\) for the intended algebraic parameter space of these algebras with Hilbert function \(T\), once its construction and topology have been specified. A nonempty topological space is irreducible if it is not the union of two proper closed subsets; its irreducible components are its maximal irreducible closed subsets. The source asks whether, for every realizable Gorenstein sequence \(T\), \(\operatorname{Gor}(T)\) is irreducible, and, when it is reducible, asks for
\[
\#\operatorname{Irr}\bigl(\operatorname{Gor}(T)\bigr),
\]
where \(\operatorname{Irr}(X)\) denotes the set of irreducible components of \(X\).

This remains a qualified formulation: a collection of algebras or their isomorphism classes alone does not specify an algebraic parameter space or its topology. The ground field, residue-field convention, and intended parameterization must be recovered from the source before the component question is fully determined. The supplied dated curation labels the problem partially solved and reports that universal irreducibility fails for known Hilbert functions, with irreducibility and component counts determined only in selected cases. Those claims are supplied assessments, not independently verified conclusions of this draft.
\subsection{Short English statement}
For a Hilbert function realized by a local Artinian Gorenstein algebra, ask whether the intended parameter space of all such algebras is irreducible and, if it is reducible, determine its number of irreducible components. The field and parameter space are unspecified. Supplied dated curation reports partial results and counterexamples to universal irreducibility; this draft requires mathematical review.
\subsection{Sources}
\begin{itemize}
\item Local UnsolvedMath registry: supplied record with permanent ID \(30004529\), title ``Irreducibility of Gorenstein-Algebra Families with Fixed Hilbert Function,'' source problem code \texttt{catalogue entry}, in the Oberwolfach Reports Open Problems collection. The proposed year is recorded as 2020; no individual proposer is supplied.
\item Supplied report citation: \emph{Lefschetz Properties in Algebra, Geometry and Combinatorics} (2021), Oberwolfach Reports, DOI 10.4171/owr/2020/31, pp.\ 1539--1600. \url{https://doi.org/10.4171/owr/2020/31}. The record describes the report as documenting workshop discussions and open problems; it gives no specific page or speaker attribution for this question.
\item Supplied literature source labeled ``Iarrobino--Macias Marques'': \url{https://arxiv.org/abs/2112.14664}.
\item Supplied literature source labeled ``Gorenstein component studies'': \url{https://arxiv.org/abs/2302.00287}.
\item Literature assessment checked at \texttt{2026-08-22}: the supplied curation labels the problem partially solved, reports reducibility for known Hilbert functions, and reports irreducibility and component counts only in selected cases. This is supplied dated curation derived from the OWR record, not independently verified literature assessment. The record expressly requests independent verification of time-sensitive status and cited claims.
\end{itemize}
\end{document}
